\documentclass[11pt]{article}
\usepackage[a4paper,margin=27mm]{geometry}
\usepackage{amsmath,amssymb,amsthm,xcolor,booktabs,array}
\newcommand{\PROVED}{\textcolor{green!40!black}{\textbf{PROVED}}}
\newcommand{\OPEN}{\textcolor{red!70!black}{\textbf{OPEN}}}
\newcommand{\AUDIT}{\textcolor{orange!70!black}{\textbf{AUDIT}}}
\newtheorem{theorem}{Theorem}
\newtheorem{proposition}{Proposition}
\newtheorem{lemma}{Lemma}
\title{The Witt--Loewy trace bridge: the correct Hilbert target (v99)}
\date{13 August 2026}

\begin{document}
\maketitle

\section{Why a faithful Hilbert completion is too strong}

Let $\mathcal A$ be the commutative test convolution algebra with involution
$f\mapsto f^\vee$, and let
\[
 E=\mathcal H_-^0
\]
be the odd nuclear representation of row (c).  A multiple spectral point
need not act semisimply on $E$: on a length-two local jet the action has
the form
\begin{equation}
             \pi(f)=\widehat f(\rho)I+widehat f'(\rho)N,
 \qquad N^2=0.                                      \tag{1}
\end{equation}

\begin{theorem}[Positive faithful completion kills jets --- \PROVED]
Suppose a positive Hilbert completion of the module in (1) satisfies
$\pi(f^\vee)=\pi(f)^*$ for every $f$.  If the involution fixes the simple
character (the critical case), then $N=0$.  Consequently a faithful
positive completion of the full odd Fr\'echet module would imply not only
criticality of every character but semisimplicity of every local spectral
block.
\end{theorem}

\begin{proof}
Choose $f$ whose value at $\rho$ is zero and whose first derivative there
is nonzero.  After multiplying by a scalar, the adjoint relation makes the
nonzero nilpotent $N$ either selfadjoint or skew-adjoint (the sign depends
only on the convention for $f^\vee$).  A selfadjoint or skew-adjoint
bounded operator is normal.  A normal nilpotent operator is zero, a
contradiction.  Equivalently, $\|N x\|^2=\langle N^*Nx,x\rangle$ and the
spectral theorem force $N=0$.
\end{proof}

The Weil form counts a zero with its algebraic multiplicity and does not
assert simplicity.  Therefore the correct Hilbert object cannot be a
faithful completion of $E$.  It must preserve the nuclear character while
semisimplifying the nilpotent jets.

\section{The source-defined Loewy object}

For a finite-length $\mathcal A$-module $M$, define its module radical
intrinsically by
\[
 \operatorname{rad}M:=\bigcap\{N\subset M:N\text{ is a maximal proper
 }\mathcal A\text{-submodule}\}.
\]
Iterate it and put
\begin{equation}
 \operatorname{gr}_{L}M
 :=\bigoplus_{j\ge0}\operatorname{rad}^jM/
                         \operatorname{rad}^{j+1}M.  \tag{2}
\end{equation}
This definition uses only the source module and its action; it contains no
zero of $\xi$ and no sign choice.

For the nuclear representation $E$, the candidate global construction is to
apply (2) to its finite-length Riesz subquotients and take a nuclear locally
convex direct limit.  Equivalently, one would apply (2) to every finite
spectral subquotient in a summability filtration.  Denote this candidate by
\begin{equation}
                      E^{\rm ss}:=\operatorname{gr}_{L}E.       \tag{3}
\end{equation}

\begin{theorem}[Finite Loewy trace preservation --- \PROVED]
For every finite-length invariant subquotient $M$ and every test function
$f$,
\begin{equation}
          \operatorname{Tr}_{\operatorname{gr}_L M}\pi(f)
          =\operatorname{Tr}_{M}\pi(f).             \tag{4}
\end{equation}
In particular the finite construction preserves every algebraic
multiplicity, even though it deletes all nilpotent couplings.
\end{theorem}

\begin{proof}
The radical filtration is invariant under
every $\pi(f)$.  In a basis adapted to the filtration, $\pi(f)$ is block
upper triangular and its diagonal blocks are precisely its actions on the
summands of (2).  The trace of a block upper triangular matrix is the sum
of the traces of its diagonal blocks, proving (4).
\end{proof}

To promote (4) unconditionally to the topological object (3), one must prove that the finite
Loewy pieces assemble in a nuclear topology and that their traces converge
to Meyer's character in its stated distributional topology.  Row (c)
proves summability of the original quotient, but does not state this Loewy
assembly.  The global promotion is therefore \AUDIT, not used below as an
unconditional lemma.

\section{Positive simple factors and the critical line}

Every simple factor of $E^{\rm ss}$ is a one-dimensional character module.
Under the subsequent spectral description of row (c), its character is
\begin{equation}
       \chi_\rho(f)=\widehat f(\rho),                \tag{5}
\end{equation}
with algebraic repetitions supplied by the different Loewy layers.
The source involution obeys
\begin{equation}
 \chi_\rho(f^\vee)
 =\overline{\widehat f(1-\overline\rho)}.            \tag{6}
\end{equation}

\begin{theorem}[Simple Witt positivity test --- \PROVED]
The character module $\chi_\rho$ admits a nonzero positive Hilbert metric
for which $\pi(f^\vee)=\pi(f)^*$ for every $f$ if and only if
\begin{equation}
                       \Re\rho=\frac12.              \tag{7}
\end{equation}
\end{theorem}

\begin{proof}
On a one-dimensional positive Hilbert space the adjoint condition is
exactly
\[
 \chi_\rho(f^\vee)=\overline{\chi_\rho(f)}
 \quad\text{for every }f.
\]
By (5)--(6), this says
$\widehat f(1-\bar\rho)=\widehat f(\rho)$ for every test function.
Mellin transforms of compactly supported smooth functions separate two
distinct points, hence $1-\bar\rho=\rho$, which is (7).  Conversely, (7)
turns (6) into the required scalar adjoint identity, and the standard
metric on $\mathbb C$ is positive.
\end{proof}

Thus, at every finite spectral level, the Loewy construction removes the
artificial simplicity requirement: if a critical spectral point has
multiplicity $m$, (2) produces $m$ copies of its character and (4) keeps
the coefficient $m$ in the trace.  The asserted global direct sum is
conditional on the assembly audit above.

\section{The finite-level Witt $C^*$-factorization}

Let $G=\mathbb Q_+^\times$ and let $C^*(G)$ be the universal group
$C^*$-algebra for the centered half-Tate generators.  The row-(b) metric
gives the algebraic involution $g^*=g^{-1}$.  The action on $E^{\rm ss}$
is initially only an action of the algebraic group algebra $\mathbb C[G]$.
Let $\mathfrak M_{\rm fin}(E)$ be the family of finite-length spectral
subquotients furnished by the row-(c) nuclear spectral theorem.  Give every
one-dimensional factor of $\operatorname{gr}_L M$ its standard Hilbert
norm; the resulting operator norm is independent of the choice of a
nonzero basis in each factor.

\begin{theorem}[$C^*$-factorization criterion --- \PROVED]
The following are equivalent.
\begin{enumerate}
\item Every finite Loewy object $\operatorname{gr}_L M$ extends to a
      $*$-representation of $C^*(G)$, uniformly over
      $M\in\mathfrak M_{\rm fin}(E)$.
\item Every simple factor of $E^{\rm ss}$ is a unitary character of $G$.
\item For every finite sum $a=\sum_g a_g g$,
\begin{equation}
 \sup_{M\in\mathfrak M_{\rm fin}(E)}
 \|\pi_{\operatorname{gr}_L M}(a)\|
 \le \|a\|_{C^*(G)}
 =\sup_{\chi\in\widehat G}\left|\sum_g a_g\chi(g)\right|.      \tag{8}
\end{equation}
\end{enumerate}
Under the row-(c) spectral theorem these conditions hold if and only if all
its simple centered characters lie on the unitary axis.
\end{theorem}

\begin{proof}
A $*$-representation of a group $C^*$-algebra sends every group element to
a unitary, proving (1)$\Rightarrow$(2).  Conversely a finite Hilbert direct
sum of the simple unitary character factors, repeated according to their
Loewy multiplicities, proves (2)$\Rightarrow$(1).  Contractivity of
$*$-representations proves (1)$\Rightarrow$(3).  Conversely (8) extends
each finite algebraic representation continuously to $C^*(G)$;
preservation of the involution makes every extension a
$*$-representation.  The last assertion is the scalar test of the
preceding theorem.
\end{proof}

This gives a new, entirely source-labelled formulation of the desired
bridge: prove the single uniform contractivity estimate (8) on finite
Loewy pieces from the row-(b) metric realization.  It requires no prior
global Hilbert topology on $E^{\rm ss}$.

\begin{corollary}[One-generator reduction --- \PROVED]
Estimate (8) is equivalent to the following pair of inequalities for every
$n\ge2$ and every $M\in\mathfrak M_{\rm fin}(E)$:
\begin{equation}
 \boxed{
 \|\pi_{\operatorname{gr}_L M}(\sqrt n\,U_n)\|\le1,
 \qquad
 \|\pi_{\operatorname{gr}_L M}(n^{-1/2}U_{1/n})\|\le1.}       \tag{8c}
\end{equation}
Thus the missing promotion can be tested one correspondence at a time.
\end{corollary}

\begin{proof}
Condition (8) gives (8c) because both centered group elements have
$C^*(G)$-norm one.  Conversely, on a simple factor let $\lambda_n$ be the
eigenvalue of $\sqrt n\,U_n$.  The second operator in (8c) is its inverse,
so the two inequalities give $|\lambda_n|\le1$ and
$|\lambda_n|^{-1}\le1$.  Hence $|\lambda_n|=1$ for every $n$.  Every
finite Loewy graded object is a direct sum of these unitary characters, so
it is a $*$-representation of $C^*(G)$.  The factorization criterion then
gives (8).
\end{proof}

Row (b) proves that $n^{-1/2}$ is the norm of its based metric contact line
and that the algebraic Witt generator has the corresponding adjoint.  It
does not prove (8c): (8c) is an operator-norm assertion on the simple
subquotients of the different completed object $E$.  A contractive functor
from the row-(b) metric correspondence category to those subquotients would
prove (8c), but the functor currently constructed in row (c) is monoidal
and trace-compatible, not norm-enriched.

\begin{theorem}[Global Hilbert assembly after (8) --- \PROVED]
If (8) holds, there is a positive Hilbert representation
\begin{equation}
 \mathscr H_-^{\rm ss}
 =\widehat\bigoplus_{\text{simple factors of }E}
        \mathbb C_{\chi},                            \tag{8a}
\end{equation}
where every factor is repeated with its Loewy multiplicity, such that the
integrated action of every test function is trace class and
\begin{equation}
 \operatorname{Tr}_{\mathscr H_-^{\rm ss}}\pi(f)=\tau_-(f).    \tag{8b}
\end{equation}
Thus the global topological audit following (4) is not an additional
hypothesis for the implication (8)$\Rightarrow$ row (d).
\end{theorem}

\begin{proof}
By the factorization criterion, every simple centered scaling character is
unitary.  Form the Hilbert direct sum (8a), retaining the number of copies
given by the finite Loewy filtration.  On the subsequent spectral
description of row (c), these characters are $u\mapsto e^{i\gamma u}$.
The Mellin transform of a compactly supported smooth test function decays
faster than every power of $|\gamma|$, while the nuclear spectral counting
used in row (c) has polynomial (indeed $O(R\log R)$) growth.  Hence the
diagonal eigenvalue sequence is absolutely summable and the integrated
operator is trace class.  Finite trace preservation (4), followed by
absolute convergence, gives (8b).  No rearrangement of a conditionally
convergent series occurs.
\end{proof}

\section{What is and is not closed}

Let $\tau_-^{\rm ss}$ be the trace on the Hilbert direct sum (8a).  The
preceding theorem shows that if (8) holds then
\[
 \tau_-^{\rm ss}(f)=\tau_-(f).
\]
For primitive $f$ the even Tate trace vanishes, and therefore
\begin{equation}
 B_{\rm nuc}(f,f)
 =-\operatorname{Tr}\bigl(\pi_{E^{\rm ss}}(f)
                  \pi_{E^{\rm ss}}(f)^*\bigr)\le0.  \tag{9}
\end{equation}
So (8) closes row (d) without requiring simplicity of the zeros and without
requiring a faithful Hilbert completion of Meyer's nonsemisimple quotient.

The files currently prove the algebraic involution on the multiplier
labels and the equality of their scalar contact weights.  They do not prove
the $C^*$-norm estimate (8) for the completed Loewy object.  Equality of
traces on individual labels is weaker than complete contractivity: a
one-dimensional algebraic character $g\mapsto g^\alpha$ has perfectly
defined label traces for any $\alpha\in\mathbb C$, but factors through
$C^*(G)$ only when $\Re\alpha=0$.

\begin{center}
\begin{tabular}{>{\raggedright\arraybackslash}p{.55\linewidth}
                >{\centering\arraybackslash}p{.16\linewidth}
                >{\raggedright\arraybackslash}p{.20\linewidth}}
\toprule
Statement & Status & Evidence \\
\midrule
Faithful completion is unnecessarily stronger than row (d) & \PROVED & (1) \\
Finite source-defined Loewy semisimplification & \PROVED & (2) \\
Finite trace and multiplicity preservation & \PROVED & (4) \\
Unconditional topological Loewy direct limit & \AUDIT & topology/trace limit not supplied \\
Global trace-class Hilbert assembly after (8) & \PROVED & (8a)--(8b) \\
Positive metric on each simple factor exactly at the unitary axis & \PROVED & (5)--(7) \\
Reduction to one $C^*$-contractivity estimate & \PROVED & (8) \\
Derivation of (8) from rows (a)--(c) & \OPEN & no such norm comparison in the files \\
Row (d), condition $G$ & \OPEN & follows from the still-open estimate (8) \\
\bottomrule
\end{tabular}
\end{center}

The mathematical gain is that the completion problem has now been typed
correctly.  The missing theorem is not ``put a Hilbert norm on Meyer'' and
not ``assume all zeros simple''.  It is the promotion of the
semisimplified Witt action from $\mathbb C[G]$ to $C^*(G)$ with its
source-fixed multiplicities.  Any proof of that promotion must produce
(8), not merely reproduce the scalar contacts $\Lambda(n)$.

\end{document}
